\documentclass[11pt, a4paper, oneside]{article}
\usepackage[ngerman]{babel}
\usepackage{worksheet}

\begin{document}
	\author{L. Bung}
	\title{Exponentialgleichungen und Schnittpunkte}
	\subject{Mathematik}
	\class{FHR1}
	\maketitle
	
	\task{Exponentialgleichungen 1}
	
	Lösen Sie die Exponentialgleichung.
	
	\begin{enumerate}[label=\alph*)]
		\item $3^x \cdot \left(3^x - 2\right) = 0$
		\item $\left(2^x - 1\right) \cdot \left(2^{-x} + 1\right) = 0$
		\item $e^{2x} - 7e^x = 0$
	\end{enumerate}
	
	\checkered[16cm]
	
	\begin{hint}{Exponentialgleichungen durch Substitution lösen}
		\phantom{x}\vspace{5cm}
	\end{hint}
	
	\task{Exponentialgleichungen 2}
	
	Lösen Sie die Exponentialgleichung.
	
	\begin{enumerate}[label=\alph*)]
		\item $\left(2^x\right)^2 - 6 \cdot 2^x + 8 = 0$
		\item $4^x - 12 \cdot 2^x = -32$
		\item $4 \cdot e^{-x} + 5 + e^x = 0$
	\end{enumerate}
	
	\checkered[11cm]
	
	\task{Schnittpunkte von Exponentialfunktionen}
	
	\begin{enumerate}[label=\alph*)]
		\item Welche Möglichkeiten gibt es für die gegenseitige Lage von zwei Exponentialfunktionen?
		Begründen Sie Ihre Antwort!
		
		\lines[2cm]
		
		\item Berechnen Sie die Schnittpunkte der beiden Exponentialfunktionen.
		Sie können Ihre Lösung mithilfe von Geogebra überprüfen.
		\begin{enumerate}[label=(\roman*)]
			\item $f(x) = 3 \cdot 2^x - 5$ und $g(x) = 2^x + 1$
			\item $f(x) = 3^x$ und $g(x) = 3 \cdot 0.5^x$
			\item $f(x) = 1.5 \cdot e^x$ und $g(x) = 3 \cdot e^{-x}$
		\end{enumerate}
	\end{enumerate}
	
	\checkered[14.5cm]
	
	\let\clearpage\relax
\end{document}